\vspace{4pt}
\noindent
\textbf{Cumulative Energy Knee–Based Subspace Partitioning.}
To determine the boundary between the dominant and residual spectral subspaces,
SALT–EDoRA employs a \emph{cumulative energy knee} criterion that adaptively selects
\(r_{\text{top}}\) from the intrinsic geometry of the singular value spectrum.
Given the ordered singular values \(\{\sigma_i\}_{i=1}^{k}\) where \(k = \min(m,n)\),
two normalised sequences are constructed:
the cumulative spectral energy
\(e_i = \bigl(\sum_{j=1}^{i} \sigma_j^2\bigr) \big/ \bigl(\sum_{j=1}^{k} \sigma_j^2\bigr)\)
and the uniform index axis
\(t_i = i / k\),
both of which lie in \([0, 1]\).
The deviation \(e_i - t_i\) measures how far the actual energy accumulation exceeds
that of a flat, uninformative spectrum at each rank.
The knee is defined as the index of maximum deviation,
\begin{equation}
    r_{\text{top}} \;=\; \operatorname*{arg\,max}_{i}\,\bigl(e_i - t_i\bigr),
    \label{eq:knee}
\end{equation}
corresponding to the \emph{elbow} of the scree plot—the rank at which adding further
singular values to the head subspace yields diminishing proportional energy returns.
Intuitively, all singular vectors to the left of this knee constitute the compact,
high-energy regime that SALT's scale-and-shift parameterisation is designed to
preserve; those to the right form the dispersed, low-energy residual that EDoRA's
intrinsic-rank reparameterisation targets.
To preclude degenerate splits in layers whose spectra are nearly flat or
pathologically sharp, the knee index is clamped to
\(r_{\text{top}} \in [\alpha_{\min} k,\, \alpha_{\max} k]\),
where \(\alpha_{\min} = 0.50\) and \(\alpha_{\max} = 0.90\) ensure
that the head always constitutes the majority of the spectral mass
while leaving a non-trivial residual subspace for the tail update.
This adaptive partitioning requires no tuning beyond the two fractional bounds
and reflects the intrinsic spectral structure of each weight matrix,
providing a principled foundation for allocating parameters between
spectral scale-shift adaptation and low-rank directional reparameterisation.
